\chapter{任务五：Agent Loop substrate 与 Workflow EEVDF}
\label{chap:task5}

\section{约束}

Agent Loop 的“循环”不应等同于 while 轮询。真正的运行时需要在等待用户、模型、文件
变化、IPC 或 heartbeat 时让线程睡眠；事件到达后按可信 route 唤醒；多个 Agent 或多个
workflow 并发时，调度器还要避免线程放大、短睡眠投机和单一高频事件长期占用 CPU。

内核也不能越界成为 planner。我们提供 watch/wait、heartbeat、MESSAGE route、LLM
correlation、lifecycle 和 workflow 调度 substrate，下一步工具仍由 Guest 策略或模型
决定。

\section{有界事件队列与 watch/wait}

每个 Agent 具有容量 16 的事件队列，其中保留 kernel reserve，并限制同 source 最多
4 条排队事件。route 表同样有界。发送 MESSAGE 需要 \code{MESSAGE_SEND} capability、同 active
workflow、显式 source/target route 与匹配 watch；内核把真实 source/target/control、
corr 和 provenance 写入事件，不相信 payload 声明的 actor。

\code{agent_wait} 使用 thread-generation reservation，先检查已排队事件，再原子安装
waiter 并睡眠，避免“检查为空后事件到达、随后永久睡眠”的 lost wakeup。timeout、
cancel 和 event consume 都更新 loop state 与 Context/audit。无事件时线程不忙轮询。

\section{Heartbeat 是事件，不是内核业务}

Agent 可以设置 heartbeat interval。timer 到期后，内核只产生 timer event 并唤醒
Agent；它不执行恢复脚本或模型逻辑。用户态收到 heartbeat 后决定检查状态、发送消息
或继续等待。停止 heartbeat 会清理 generation-bound 状态，旧 timer 不能命中新进程。

这个机制覆盖长期会话的 liveness，但不提供分布式 failure detector 的准确性保证。
heartbeat 延迟仍受可调度 safe point 和当前单 Hart 负载影响。

\section{绑定身份的 LLM correlation RPC}

\code{LLM_REQUEST/LLM_RESPONSE} 复用 MESSAGE 和 wait substrate。请求成功后，内核登记
完整 lifecycle、requester PID/control、relay PID/control、corr 和 deadline。corr 必须
非零并对同 requester 身份严格递增；只有 MESSAGE 真正入队且 pending 发布成功才推进，
失败投递不消耗编号。

只有原 Relay 在同 lifecycle 返回匹配 corr，才产生 \code{LLM_DONE} 并原子消费 pending。
pending 使用 120 秒 kernel-tick TTL；全局容量为 NPROC，每 requester 上限为队列容量。
有界 terminal history 在记录保留时区分 consumed replay 的 \code{STALE} 与 expired 的
\code{TIMEOUT}；覆盖后旧响应按 unmatched \code{DENIED}。该 RPC 不解析 prompt、JSON、
HTTPS 或 MCP。

\section{Workflow Credit Domain}

资源账户维护 used(U)、pending(P)、free(F) 和 \code{held=U+P+F}。额度不足时按资源种类
quantum 从全局预充少量 F，热路径只执行：

\begin{center}
\code{reserve: F→P}\quad
\code{commit: P→U}\quad
\code{cancel: P→F}\quad
\code{release: U→F}
\end{center}

这些本地移动不改变 held，减少全局统计写入；补充 credit 前仍检查 global、class 和
account hard limit。压力路径、context switch、account close 和 fence 会 trim F。
fence 在单锁中 \code{trim+snapshot}，要求全部 P 为零，导出精确 U 和 account
slot/generation、credit epoch、U vector 的 digest。F 是预充额度，不是超卖。

\section{按 workflow 聚合的 EEVDF}

传统按线程公平会奖励创建更多线程的 workflow。我们把同一 resource domain 聚合为
一个调度实体。调度器用 workflow vruntime 计算 lag，只有
\code{vruntime <= global_vtime} 的实体 eligible；latency class 和 event/deadline 缩短
service request，再从 eligible 集合选择最早 virtual deadline。dispatch 后按实际
service cycles 给 workflow 记账，睡眠 lag 向 global vtime 衰减。

一个实体走 fast path。身份不完整、内部不变量失败或容量外情况回退旧 RR/Agent
heuristic，并计数 fallback；失败计划不修改 EEVDF 状态。实体表总容量为 4，其中固定
包含一个 \code{BOOT_SEALED} bootstrap participant，因此最多同时有三个 fresh workflow。

所谓 1/4/16 场景必须精确解释：1 是单实体 fast path；4 是 bootstrap 加三个 fresh 的
满容量；16 是四波逻辑样本，每波复用同一 bootstrap 并累计创建 12 个 fresh lifecycle，
不是 16-way 并发，也不是每波四个 fresh。

\section{调度实测}

one-shot campaign 用六个 fresh QEMU boot 采集并发度 1、2、3、4 的 service window 与
exact wake probe。504 条 probe 直接读取 scheduler trace 的 \code{ready_age}，其中
425 条为 0 tick、79 条为 1 tick，没有右删失。tick 为 10 ms；这说明在该负载下全部
fresh waiter 在一个 tick 内被观察到调度，不等于微秒级 wake latency。

\ReportFigure[0.94\textwidth]
  {../figures/performance/05_eevdf_latency_fairness.pdf}
  {../figures/performance/05_eevdf_latency_fairness.pdf}
  {EEVDF exact wake latency 与并发 1--4 的 Jain fairness}
  {fig:eevdf-performance}

以同一窗口 raw service cycles 重算的 Jain 中位数在并发 1/2/3/4 下分别为
1.000000、0.999985、0.999993、0.999985。并发 1 的 Jain=1 是平凡基线；图使用 focused
y-axis 展开接近 1 的差异。公平数值只针对已采集 workflow service，不等于整机所有
进程或跨硬件公平性。

EEVDF boot 05 有三条并发串口 sample summary 拼接。提取器保留 warning 并排除不完整
summary，504 条带完整前缀的 exact wake probe 全部保留；我们没有静默修补或插值。

\section{实现与验证}

\code{os/agent_ipc.c} 实现队列、watch/wait、route、heartbeat 和唤醒；
\code{os/agent_core.c} 实现 LLM pending 与 Agent loop state；
\code{os/workflow_scheduler.c} 实现 EEVDF；resource controller 与 credit domain 维护账户。
\code{agentloop_ucore} 覆盖 atomic wait/wakeup、timeout、cancel 和 heartbeat；
\code{agentsched_ucore}、\code{agent_eevdf_ucore} 覆盖 fast path、满容量、公平、线程放大、
sleep decay、fallback 和 wake histogram。Host scheduler model 单独验证抽象 lag/deadline
算法，但不代替 Guest 固定槽位和真实 lifecycle 拓扑。

\begin{evidencebox}[任务五结论]
我们实现了事件驱动、无事件休眠、可信多 Agent IPC、模型 correlation 和 workflow
级调度。内核提供 Loop substrate，不解释目标或选择下一工具；hard deadline 在首个
schedulable safe point 结算，不承诺 wall-clock 终止上界。
\end{evidencebox}
